\begin{graphillustration}{Путь $P_4$: среднее ребро — паросочетание, максимальное по включению (ни одно ребро не добавить), но не наибольшее. Два крайних ребра — совершенное паросочетание.}{fig:notation-matchings}
\begin{tikzpicture}[font=\small]
\begin{scope}
\node[gi free] (a) at (0,0) {$a$};
\node[gi vertex] (b) at (1.3,0) {$b$};
\node[gi vertex] (c) at (2.6,0) {$c$};
\node[gi free] (d) at (3.9,0) {$d$};
\draw[gi edge] (a) -- (b);
\draw[gi match] (b) -- (c);
\draw[gi edge] (c) -- (d);
\node[gi note] at (1.95,-0.8) {$M=\{bc\}$: максимально по включению,\\$a,b,c,d$ — дополняющий путь};
\end{scope}
\begin{scope}[shift={(6.5,0)}]
\node[gi vertex] (a) at (0,0) {$a$};
\node[gi vertex] (b) at (1.3,0) {$b$};
\node[gi vertex] (c) at (2.6,0) {$c$};
\node[gi vertex] (d) at (3.9,0) {$d$};
\draw[gi match] (a) -- (b);
\draw[gi edge] (b) -- (c);
\draw[gi match] (c) -- (d);
\node[gi note] at (1.95,-0.8) {$M'=\{ab,cd\}$: совершенное,\\$|M'|=\alpha'(P_4)=2$};
\end{scope}
\end{tikzpicture}
\end{graphillustration}
